\begin{frame}{Caso 2. Rotación de cuentas por cobrar}
\smallcase
\begin{columns}[T,totalwidth=\textwidth]
\column{0.58\textwidth}
\casehead{Enunciado.} Servicios del Pacífico S.A.C. desea determinar la rapidez con la que convierte ventas a crédito en efectivo.

\casehead{Datos.} Ventas a crédito = \soles{1,200,000.00}; CxC inicial = \soles{180,000.00}; CxC final = \soles{220,000.00}.

\casehead{Operación.}
\resizebox{\linewidth}{!}{%
\begin{tabular}{@{}ll@{}}
Promedio CxC: & (\soles{180,000.00} + \soles{220,000.00}) / 2 = \soles{200,000.00}\\
Rotación: & \soles{1,200,000.00} / \soles{200,000.00} = 6.00 veces\\
Periodo: & 360 / 6.00 = 60.00 días
\end{tabular}}
\casehead{Respuesta.} Rotación = 6.00 veces; periodo promedio = 60.00 días.

\casehead{Interpretación.} La empresa cobra aproximadamente seis veces al año y tarda 60 días en recuperar ventas a crédito.

\column{0.37\textwidth}
\centering
\begin{tikzpicture}[node distance=0.45cm, every node/.style={font=\scriptsize, draw, rounded corners=2pt, align=center, minimum width=2.2cm}]
\node (a) {Venta a\\crédito};
\node[right=of a] (b) {Cuenta por\\cobrar};
\node[below=of b] (c) {Cobranza};
\node[left=of c] (d) {Efectivo};
\draw[-{Stealth}] (a) -- (b);
\draw[-{Stealth}] (b) -- (c);
\draw[-{Stealth}] (c) -- (d);
\draw[-{Stealth}] (d) -- (a);
\end{tikzpicture}
\end{columns}
\end{frame}

\begin{frame}{Caso 3. Rotación de inventarios}
\smallcase
\begin{columns}[T,totalwidth=\textwidth]
\column{0.58\textwidth}
\casehead{Enunciado.} Una empresa comercial desea conocer cuántas veces renueva su inventario y cuánto permanece la mercadería en almacén.

\casehead{Datos.} Costo de ventas = \soles{900,000.00}; inventario inicial = \soles{140,000.00}; inventario final = \soles{160,000.00}.

\casehead{Operación.}
\resizebox{\linewidth}{!}{%
\begin{tabular}{@{}ll@{}}
Inventario promedio: & (\soles{140,000.00} + \soles{160,000.00}) / 2 = \soles{150,000.00}\\
Rotación: & \soles{900,000.00} / \soles{150,000.00} = 6.00 veces\\
Permanencia: & 360 / 6.00 = 60.00 días
\end{tabular}}
\casehead{Respuesta.} Rotación = 6.00 veces; permanencia = 60.00 días.

\casehead{Interpretación.} El inventario se renueva seis veces al año y permanece cerca de 60 días en almacén.

\column{0.37\textwidth}
\centering
\begin{tikzpicture}[node distance=0.45cm, every node/.style={font=\scriptsize, draw, rounded corners=2pt, align=center, minimum width=2.05cm}]
\node (a) {Compra};
\node[right=of a] (b) {Almacén};
\node[below=of b] (c) {Venta};
\node[left=of c] (d) {Reposición};
\draw[-{Stealth}] (a) -- (b);
\draw[-{Stealth}] (b) -- (c);
\draw[-{Stealth}] (c) -- (d);
\draw[-{Stealth}] (d) -- (a);
\end{tikzpicture}
\end{columns}
\end{frame}
